% Part: intuitionistic-logic
% Chapter: tableaux
% Section: introduction

\documentclass[../../../include/open-logic-section]{subfiles}

\begin{document}

\olfileid{int}{tab}{int}

\olsection{引言}

分析表是带号公式的某种（向下分支的）树，即真值符（$\True$ 或
$\False$）与一个语句
\[
\sFmla{\True}{!A} \text{ 或 } \sFmla{\False}{!A}.
\]
组成的对。分析表以若干\emph{假设}开始。
每个后续的带号公式都是通过应用某条推理规则生成的。
有些推理规则在树的某个末端添加一个或多个带号公式；另一些规则则添加两个新的末端，从而产生两个分支。
应用规则所得到的带号公式，其公式部分的复杂度低于原带号公式中的公式。
当一个分支同时包含 $\sFmla{\True}{!A}$ 和
$\sFmla{\False}{!A}$ 时，我们称该分支是\emph{封闭的}。
如果一个分析表中的每个分支都是封闭的，那么整个分析表就是封闭的。
一个封闭的分析表构成一个推导，表明作为该分析表起始假设的那组带号公式是不可满足的。
这可以用来定义 $\Proves$ 关系：
$\Gamma \Proves !A$ 当且仅当存在某个有限集~$\Gamma_0 = \{!B_1,
\dots, !B_n\} \subseteq \Gamma$，使得存在一个以
\[
\{\sFmla{\False}{!A}, \sFmla{\True}{!B_1}, \dots, \sFmla{\True}{!B_n}\}.
\]
为假设的封闭分析表。

对于直觉主义逻辑，我们必须既扩展带号公式的概念，又调整联结词的规则。
在模态分析表中，公式除了带有真值符（$\True$ 或 $\False$）外，还有一个\emph{前缀}~$\sigma$。
前缀是非空的正整数序列，即 $\sigma \in
(\PosInt)^* \setminus \{\emptyseq\}$。
我们书写这些前缀时，会省略外围的 $\tuple{\ }$，并用~$.$ 而非 $,$ 来分隔各个元素。
如果 $\sigma$ 是一个前缀，那么 $\sigma.n$ 就是 $\sigma \concat \tuple{n}$；
例如，如果 $\sigma = 1.2.1$，那么 $\sigma.3$ 就是 $1.2.1.3$。
例如，
\[
\sFmla{\True}{!A \lif (!B \lif !C)}[1.2]
\]
就是一个\emph{带前缀的带号公式}（或简称为\emph{带前缀公式}）。

直观上，前缀命名了模型中某个可能满足分析表分支上公式的世界；如果 $\sigma$ 命名了某个世界，
那么 $\sigma.n$ 命名的就是一个从 $\sigma$ 所命名的世界可达的世界。

在直觉主义模型中，可达关系是自反且传递的。
就前缀而言，这意味着 $\sigma$ 从其自身可达，且任何扩展了~$\sigma$ 的前缀从 $\sigma$ 也是可达的，
即任何形如 $\sigma.n_1.\cdots.n_k$ 的前缀。我们引入记号 $\sigma.*$ 来表示 $\sigma$ 自身及其任何扩展。
换言之，前缀 $\sigma.*$ 恰好是所有从~$\sigma$ 可达的前缀。

\end{document}